\chapter[Discussão]{Discussão}
\label{cap:discussao}

Este capítulo interpreta os resultados obtidos, discute limitações do trabalho e
aponta ameaças à validade dos achados.

\section{Análise dos resultados}

O resultado central é a \textbf{ausência de vantagem \emph{consistente} do
pré-treinamento clínico}. Para a extração de relações no SemClinBr, no espaço de
três classes, o BioBERTpt não produz ganho estável sobre o BERTimbau, nem mesmo na
classe \texttt{negation\_of}, em que a especialização de domínio seria mais
plausível. Convém precisar o que os testes sustentam, porque o quadro é menos
confortável que um empate. Na métrica-alvo, o veredito de \emph{empate} vale para
as duas sementes, uma vez que o \emph{bootstrap} pareado sobre o F1 de
\texttt{negation\_of} não rejeita a hipótese nula nem na semente $42$ nem na $43$.
O que \emph{não} se mantém entre elas é a \emph{direção nessa classe}, dado que o
sinal da diferença se inverte e passa a favorecer o \emph{encoder} clínico. O teste
de McNemar, que incide sobre o conjunto dos candidatos, tampouco rejeita a hipótese
nula em qualquer das sementes, e a contagem de discordâncias aponta, nas duas, para
o \emph{encoder} clínico, por margem compatível com o acaso ($540$ contra $516$ na
semente $42$; $559$ contra $510$ na $43$), assim como a acurácia. Nenhum dos testes
aplicados detecta, portanto, diferença entre os dois \emph{encoders}. O que se
mantém, além dessa ausência de diferença, é a \emph{magnitude relativa}. A variação
do \emph{encoder} clínico entre sementes ($0{,}043$ na métrica-alvo) supera a
diferença medida entre os dois \emph{encoders} em qualquer das sementes ($0{,}015$
na $42$ e $0{,}020$ na $43$), ao passo que a variação do \emph{encoder} geral
($0{,}007$) fica abaixo dela (Seção~\ref{sec:robustez}). É nesse sentido, e apenas
nesse, que não há ganho atribuível ao domínio de pré-treino. O motivo não é uma derrota do clínico,
que chega a ficar à frente na negação sob a semente $43$, mas o fato de a diferença
entre os dois na classe investigada não sobreviver à troca de uma semente. Esse
achado contrasta com a expectativa difundida, bem como com os ganhos relatados
para o reconhecimento de entidades nomeadas clínicas em português \cite{schneider_2020_biobertpt}, de que
modelos de domínio sempre se transfiram melhor.

Duas explicações são plausíveis e compatíveis com os dados.

A primeira é \textbf{linguística}. A negação clínica em português é fortemente
lexicalizada (``não'', ``sem'', ``nega'', ``ausência de'', além de adjetivos que
embutem negação morfologicamente, como ``afebril'' e ``anictérico'', conforme o
Apêndice~\ref{ap:lexico}). Esse sinal é frequente e regular o bastante para que um
\emph{encoder} de domínio geral, com vocabulário amplo do português, o capture sem
desvantagem consistente em relação a um clínico. A caracterização do corpus reforça
essa leitura, pois a mediana da distância entre as entidades de uma relação
\texttt{negation\_of} é de um único caractere (Seção~\ref{sec:corpus}), ou seja,
o gatilho de negação está praticamente
colado ao achado negado. Detectá-lo depende pouco de conhecimento clínico
especializado.

A segunda é \textbf{estrutural}. A tarefa, tal como formulada, depende mais da
estrutura local do par, isto é, da direção, da proximidade e da delimitação das
entidades, do que de conhecimento lexical especializado de domínio, e essa
estrutura é fornecida
igualmente aos dois \emph{baselines} pela representação de entrada compartilhada
(Seção~\ref{sec:markers}). A semelhança entre as matrizes de confusão
(Seção~\ref{sec:confusao}) é a evidência mais direta disso. Os dois modelos não
apenas chegam a números próximos, mas concentram o erro no mesmo lugar, a
fronteira entre \texttt{associated\_with} e \texttt{no\_relation}.

Um corolário prático, e diretamente relevante para o desenvolvimento do
RECLin-PT, é a \textbf{eficiência}. O BERTimbau acompanha o BioBERTpt em macro-F1
médio, ficando $0{,}002$ atrás ($0{,}701$ contra $0{,}703$), com aproximadamente
$60\%$ dos parâmetros
($108{,}9$ contra $177{,}9$ milhões). Quando o desempenho é equivalente, o modelo
menor é preferível por custo de inferência e de memória, consideração concreta para
implantação em ambientes clínicos com recursos limitados. A esse argumento soma-se a
estabilidade, dado que a amplitude entre sementes do BERTimbau na métrica-alvo ($0{,}007$) é
cerca de seis vezes menor que a do BioBERTpt ($0{,}043$), o que torna seu desempenho
mais previsível em uma nova execução. Vale registrar a contrapartida. É justamente
essa dispersão que permite ao BioBERTpt superar o geral na semente $43$, de modo que
a preferência pelo \emph{encoder} menor se justifica por previsibilidade e custo,
não por superioridade média na métrica-alvo.

\section{O ganho do filtro de pista lexical}
\label{sec:discussao_filtro}

A primeira das duas explicações oferecidas na análise anterior para a ausência de
vantagem do \emph{encoder} clínico era a lexicalização forte da negação em
português. Até a comparação entre \emph{baselines}, ela permanecia justificativa
linguística,
apoiada na regularidade dos gatilhos e na mediana de distância de um único
caractere entre a pista e o achado negado. O filtro de pista a converte em medida.
Um léxico de apenas $11$ formas, induzido automaticamente do conjunto de treino,
cobre $0{,}9467$ dos pares \texttt{negation\_of} do conjunto de desenvolvimento, e
exigir a presença de uma dessas formas no primeiro argumento eleva a precisão da
classe no teste de $0{,}5547$--$0{,}6167$ para $0{,}7514$--$0{,}8059$, ao custo de
cerca de dois pontos centesimais de \emph{recall}
(Seção~\ref{sec:fase2_resultados}). O que esse resultado mostra é que a maior parte
dos falsos positivos do modelo contextual ocorre em pares que nenhum gatilho de
negação inicia, ou seja, fora do espaço em que a classe poderia legitimamente
existir. Confirma-se, assim, por medida direta e não apenas por argumento, que o
sinal de negação no português clínico é regular o bastante para ser delimitado por
uma lista curta, e é essa mesma propriedade que torna dispensável o
pré-treinamento de domínio para capturá-lo.

O erro que resta depois do filtro foi auditado, e a sua composição qualifica o
quanto ainda há a ganhar. A auditoria incide sobre uma execução sorteada entre as
quatro, com a semente do sorteio fixada em código antes do primeiro sorteio, e
recaiu sobre o RECLin-PT aplicado ao BioBERTpt na semente $43$. Essa execução
produz $33$ falsos positivos, classificados por um critério estrutural fixado antes
da leitura dos casos, que decide pela evidência do \emph{gold} e não por leitura
clínica caso a caso. Dos $33$, $18$ ($54{,}5\%$) são prováveis erros de anotação do
\emph{gold}, em que a pista e o alvo são adjacentes e o corpus não registra relação
alguma para aquela pista; $6$ ($18{,}2\%$) são ambiguidades genuínas, tipicamente
listas coordenadas\footnote{Lista coordenada é a enumeração de
achados ligados por vírgula ou conjunção e precedidos por uma única pista de
negação, como em ``sem febre, tosse ou dispneia'', em que a pista nega cada um
dos itens.} em que o \emph{gold} anota apenas o primeiro alvo; e $9$
($27{,}3\%$) são erros do modelo. Se toda a divergência de anotação fosse resolvida
a favor do modelo, os $24$ casos que não configuram erro de leitura clínica
deixariam de contar como falsos positivos, e a precisão dessa execução passaria de
$0{,}8059$ para $0{,}9384$, com F1 de $0{,}9195$. \textbf{Esse valor não é um
resultado do sistema.} Ele é um \emph{teto hipotético}, condicionado a uma
resolução favorável de divergências que esta auditoria não estabelece, e presta-se
apenas a dimensionar quanta precisão ainda resta disponível antes de se investir
GPU nessa direção. O desempenho efetivamente medido dessa execução continua sendo o
F1 de $0{,}8509$ da Tabela~\ref{tab:fase2_resultados}. Cabe registrar também o
limite da própria auditoria, cujo critério mecânico acerta o padrão dominante, que
é a lista coordenada, e erra em casos que exigiriam interpretar o que a frase quer
dizer. Uma revisão humana pode mover casos entre as três classes, e o trecho de
contexto de cada caso consta do relatório justamente para permiti-la.

\section{A classe \texttt{associated\_with}}

A classe \texttt{associated\_with} permanece a mais difícil para os dois modelos
(F1 $\approx 0{,}44$), com alta cobertura e baixa precisão. Duas causas se somam.

A primeira é \textbf{conceitual}. Distinguir associação genuína de mera proximidade
textual é intrinsecamente ambíguo em texto clínico telegráfico, em que entidades
vizinhas frequentemente não mantêm relação anotada, como ilustra o próprio
exemplo de geração de candidatos (Quadro~\ref{quad:exemplo_candidatos}), em que
cinco dos seis pares ordenados possíveis são \texttt{no\_relation}.

A segunda é \textbf{metodológica}. A janela \texttt{max\_gap}$=25$ descarta
$9{,}11\%$ das relações \texttt{associated\_with} do conjunto de teste
(Seção~\ref{sec:teto_recall}), enquanto não descarta \emph{nenhuma} relação
\texttt{negation\_of}. O F1 dessa classe é, portanto, medido sob um teto de
\emph{recall} substancialmente mais baixo que o das demais. Como o teto é idêntico
para os dois modelos, ele não afeta a comparação, objeto deste trabalho, mas
qualquer leitura do valor absoluto de F1 de \texttt{associated\_with} precisa levá-lo
em conta.

\section{Limitações}

A limitação mais séria recai sobre o alcance da caracterização de empate. O
veredito de ausência de diferença vale para as duas sementes, mas resulta de uma
\textbf{não rejeição} da hipótese nula, e não de uma demonstração de equivalência.
Soma-se a isso que a comparação pareada refeita na semente $43$ inverte o sinal da
diferença observado na semente de referência, passando a apontar para o BioBERTpt
clínico. A conclusão prática, a ausência de ganho \emph{confiável}
atribuível ao pré-treinamento clínico, sobrevive a essa inversão, porque se apoia na
instabilidade da diferença e não em sua direção; mas qualquer afirmação sobre
\emph{qual} \emph{encoder} é superior depende inteiramente da semente e não está
sustentada por estes experimentos. Além disso, \textbf{duas sementes por modelo não
bastam} para estimar um intervalo de confiança sobre essa variação, e o desvio
reportado na Tabela~\ref{tab:robustez} deve ser lido como dispersão observada, não
como estimativa de variância. Fechar a questão exige um número maior de repetições
com agregação formal.

A tarefa pressupõe \textbf{entidades já anotadas} (\emph{gold}). A avaliação não
considera o erro propagado de uma etapa de reconhecimento de entidades nomeadas
automática, de modo que os resultados refletem a extração de relações isoladamente.
Em um cenário fim a fim, os valores absolutos seriam necessariamente menores.

Por fim, o achado é específico desta \textbf{formulação de três classes} e desta
\textbf{escala de modelo} (\emph{base}). Não se pode excluir que o pré-treinamento
clínico volte a importar em espaços de relação mais ricos, em modelos maiores ou em
cenários com menos dados de ajuste.

\section{Ameaças à validade}

Quanto à \textbf{validade interna}, o congelamento das partições com semente fixa e
verificação por \emph{hashes} SHA-256 (Apêndice~\ref{ap:sha256}), somado ao
particionamento em nível de documento, mitiga o vazamento de dados, a principal
causa de resultados não reproduzíveis em PLN biomédico. A paridade total de
implementação entre os dois \emph{baselines} (Capítulo~\ref{cap:prototipo}) garante
que a diferença medida reflita o \emph{checkpoint}, e não artefatos de código.
Restam, contudo, decisões de projeto que influenciam os números absolutos, das
quais a mais consequente é a janela \texttt{max\_gap}$=25$ na geração de pares
candidatos. Esse limite foi quantificado por partição e por classe
(Seção~\ref{sec:teto_recall}) e, portanto, é uma limitação \emph{documentada e de
magnitude conhecida}, e não uma fonte de incerteza não medida.

Uma segunda ameaça diz respeito à \textbf{independência das instâncias}. O
\emph{bootstrap} pareado trata os $19\,210$ candidatos como independentes, o que
não é estritamente verdadeiro, pois pares de um mesmo documento compartilham texto e
vocabulário, de modo que o intervalo de confiança reportado é provavelmente mais
estreito que o real. Essa ameaça atenua-se, mas não desaparece, pelo fato de a
comparação ser \emph{pareada}, já que ambos os modelos são avaliados exatamente sobre as
mesmas instâncias, e a dependência afeta os dois igualmente. Ainda assim, a
conclusão de empate nas duas sementes deve ser lida com essa ressalva.

Quanto à \textbf{validade externa}, o SemClinBr, apesar de multi-institucional e
multiespecialidade, representa um conjunto específico de instituições e estilos de
redação; a generalização para outros contextos clínicos e para outros espaços de
relação deve ser feita com cautela. As inconsistências de anotação conhecidas do
corpus, ainda que tratadas e registradas pelo \emph{parser}, também impõem um limite
superior ao desempenho alcançável, limite que se aplica igualmente aos dois
modelos comparados.

\subsection{Deriva dos \emph{offsets} de entidade}
\label{sec:ameaca_offsets}

Uma última ameaça à validade interna foi identificada durante a construção do
RECLin-PT e merece registro próprio, por não ter sido considerada no protocolo
original. Os \emph{offsets} das entidades do corpus não indexam corretamente o
texto do documento na representação canônica. No conjunto de teste, apenas
$64{,}4\%$ das entidades satisfazem a identidade entre o trecho recortado pelo
\emph{offset} e o texto anotado da própria entidade, contra $58{,}9\%$ no treino e
$65{,}8\%$ no desenvolvimento. A deriva é cumulativa e negativa ao longo do
documento, assinatura compatível com uma normalização de espaço em branco aplicada
ao texto \emph{depois} do cálculo dos \emph{offsets}. A etapa exata em que ela
ocorre permanece por identificar, e a correção automática de deslocamentos de
$\pm 1$ e $\pm 2$ caracteres executada pelo \emph{parser}
(Capítulo~\ref{cap:metodologia}) não alcança uma deriva acumulada dessa magnitude.

Três razões circunscrevem o alcance dessa ameaça. A primeira é que ela afeta os
dois \emph{encoders} de forma idêntica, por atuar na construção da janela marcada
(Seção~\ref{sec:markers}), que é compartilhada, de modo que a comparação do
Capítulo~\ref{cap:experimentos} permanece válida. A segunda é que o filtro de pista
do RECLin-PT não consulta \emph{offset} algum, pois decide pela forma de superfície
registrada na própria entidade, e o resultado da Seção~\ref{sec:fase2_resultados}
não depende, portanto, desse alinhamento. A terceira é que a distância entre
entidades empregada na geração de candidatos e na regra pura é calculada entre
\emph{offsets} mutuamente consistentes, permanecendo válida como medida relativa.

O que a deriva sugere é um teto de desempenho a recuperar, e não um erro de
comparação. Em cerca de $36\%$ das entidades do conjunto de teste, os marcadores
\texttt{[E1]} e \texttt{[E2]} são inseridos alguns caracteres deslocados em relação
à superfície que deveriam delimitar, e todos os valores absolutos deste trabalho
foram medidos sob essa condição. Corrigir o alinhamento é, por isso, uma das
continuidades previstas no Capítulo~\ref{cap:conclusao}.
